\hypertarget{priest-flyxion-correspondence-ledger-rev.-4}{%
\chapter*{Priest-Flyxion Correspondence Ledger (rev. 4)}
\addcontentsline{toc}{chapter}{Priest-Flyxion Correspondence Ledger (rev. 4)}
\markboth{Priest-Flyxion Correspondence Ledger (rev. 4)}{Priest-Flyxion Correspondence Ledger (rev. 4)}\label{priest-flyxion-correspondence-ledger-rev.-4}}

The matrix's function changes here. It no longer records what looked
formally alike before drafting; it records what the finished six-part
draft actually established, chapter by chapter, against the manuscript
files themselves (all six parts are drafted and were checked directly for
this revision, not reconstructed from progress summaries). Four columns:
the proposed correspondence, where it was tested, its final status, and
the exact qualification or open obligation attached to that status.

\textbf{Status vocabulary, revised (rev. 4) to separate three things ``Bridge''
was doing at once:}
- \textbf{Direct --- established}: the chapter actually supplied the
translation, not just a resemblance.
- \textbf{Bridge}: the manuscript develops a substantive connection but
stops short of derivation, without pinning down one specific missing
step --- the gap is real but not yet nameable as a single result.
- \textbf{Open obligation}: the manuscript develops the connection far enough
to expose one specific, nameable missing result --- a proof, a
derivation, or a technical question stated precisely enough that
answering it would close the gap. Stronger than Bridge because the
debt is specific, not diffuse.
- \textbf{Tension --- retained}: a genuine disagreement, kept rather than
resolved, and not to be papered over later.
- \textbf{Rejected}: a plausible-sounding correspondence was investigated and
found to conflate two genuinely different structures. Stronger than
``None'' --- this is a negative result, not an absence of one.
- \textbf{Distinct but related}: two structures address the same territory
without being the same object; neither Direct nor Rejected fits.
- \textbf{None}: no useful correspondence found; rare enough not to need its
own section.
- \textbf{Unpursued} (excluded from the tables below, listed separately at
the end): proposed during planning, but no chapter ever tested it.
This is not a status the finished manuscript can be scored against ---
the manuscript simply doesn't contain an attempt. Keeping these in the
main tables alongside Bridge and Open obligation would misstate what
the ledger claims to record: status \emph{after drafting}, for material
that was actually drafted.

\begin{center}\rule{0.5\linewidth}{0.5pt}\end{center}

\hypertarget{part-i-logic-as-constraint}{%
\section{Part I --- Logic as Constraint}\label{part-i-logic-as-constraint}}

\begin{longtable}[]{@{}
  >{\raggedright\arraybackslash}p{(\columnwidth - 6\tabcolsep) * \real{0.2500}}
  >{\raggedright\arraybackslash}p{(\columnwidth - 6\tabcolsep) * \real{0.2500}}
  >{\raggedright\arraybackslash}p{(\columnwidth - 6\tabcolsep) * \real{0.2500}}
  >{\raggedright\arraybackslash}p{(\columnwidth - 6\tabcolsep) * \real{0.2500}}@{}}
\toprule\noalign{}
\begin{minipage}[b]{\linewidth}\raggedright
Correspondence
\end{minipage} & \begin{minipage}[b]{\linewidth}\raggedright
Tested in
\end{minipage} & \begin{minipage}[b]{\linewidth}\raggedright
Status
\end{minipage} & \begin{minipage}[b]{\linewidth}\raggedright
Qualification / open obligation
\end{minipage} \\
\midrule\noalign{}
\endhead
\bottomrule\noalign{}
\endlastfoot
Classical consequence ↔ admissibility at maximal strictness & Ch.1 & Direct --- established & Classical logic is R turned to maximal strictness in every direction, not a default. \\
Logical consequence as truth-preservation ↔ \(x \leadsto_R y\) & Ch.1 & Direct --- established & Same shape as the ordinary definition, with admissibility as the varying parameter. \\
Explosion ↔ unrestricted propagation closure & Ch.1, Ch.4 & Direct --- established & Explosion is a specific, substantive closure choice, not a logical necessity. \\
Lewis's argument for explosion ↔ repair cost formulation & Ch.1 & Open obligation & Named missing result: a formal derivation from \(E(a)+\lambda R(a)\), not yet supplied --- currently gestured at only. \\
\end{longtable}

\hypertarget{part-ii-worlds-and-admissibility}{%
\section{Part II --- Worlds and Admissibility}\label{part-ii-worlds-and-admissibility}}

\begin{longtable}[]{@{}
  >{\raggedright\arraybackslash}p{(\columnwidth - 6\tabcolsep) * \real{0.2500}}
  >{\raggedright\arraybackslash}p{(\columnwidth - 6\tabcolsep) * \real{0.2500}}
  >{\raggedright\arraybackslash}p{(\columnwidth - 6\tabcolsep) * \real{0.2500}}
  >{\raggedright\arraybackslash}p{(\columnwidth - 6\tabcolsep) * \real{0.2500}}@{}}
\toprule\noalign{}
\begin{minipage}[b]{\linewidth}\raggedright
Correspondence
\end{minipage} & \begin{minipage}[b]{\linewidth}\raggedright
Tested in
\end{minipage} & \begin{minipage}[b]{\linewidth}\raggedright
Status
\end{minipage} & \begin{minipage}[b]{\linewidth}\raggedright
Qualification / open obligation
\end{minipage} \\
\midrule\noalign{}
\endhead
\bottomrule\noalign{}
\endlastfoot
Possible worlds ↔ admissible states \(w \in \operatorname{Adm}(R)\) & Ch.5 & Direct --- established & Deliberately ontology-neutral; realism/actualism/Meinongian debate left open on purpose. \\
Kripke necessity ↔ single-parameter invariance & Ch.5 (provisional), Ch.6 (corrected) & \textbf{Rejected as originally framed}, replaced & The single-\(R\) box formula conflates admissibility and accessibility; Ch.6 splits them into \((R, \mathcal{R})\). The rejection is of the \emph{conflation}, not of necessity-as-invariance generally. \\
Normal modal systems (K, T, S4, S5) ↔ structural constraints on \(\mathcal{R}\) & Ch.6 & Direct --- established & Reflexivity/transitivity/symmetry imposed on \(\mathcal{R}\) specifically, leaving \(R\) untouched. \\
Non-normal worlds ↔ refusal & Ch.7 & \textbf{Rejected}, replaced by non-normal worlds ↔ state-indexed regime \(R_w\) & Refusal is a targeted, Level-3 decline of one transformation; non-normal worlds suspend the compositional evaluation rule itself --- a different structural level. The original Direct rating was vocabulary-matching, not a derivation. \\
Impossible worlds ↔ locally varying admissibility & Ch.7 & Direct --- established, with a downstream open obligation & \(R_w\) formalizes this cleanly at the level of single states; the obligation surfaces one part later (see Part V/VI section below), not here. \\
Inconsistency / impossibility / non-normality as one phenomenon & Ch.7, using Part III as control case & Rejected as a collapse; three independent axes established & ``Impossibility is indexed to a constraint regime; contradiction is indexed to support'' --- stated as this part's foundational sentence. \\
Strict implication, conditional logics ↔ admissibility transport \(p \Rightarrow_R q\) & Ch.8 & Open obligation & Named missing result: the transformation-selection rule Lewis's similarity ordering supplied has no replacement yet. \\
Distinguishability geometry ↔ Lewisian similarity spheres & Ch.8 & Open obligation & Chapter ends on a precisely stated question (can similarity be derived from admissibility geometry) rather than claiming the replacement --- a nameable research question, not a diffuse gap. \\
\end{longtable}

\hypertarget{part-iii-gaps-gluts-and-propagation}{%
\section{Part III --- Gaps, Gluts, and Propagation}\label{part-iii-gaps-gluts-and-propagation}}

*(Editorial note: the working title ``Gaps, Gluts, and Obstructions'' is
retired for this part. Part VI earned the right to use ``obstruction'' as
a technical term distinct from contradiction; using the same word in
Part III's own heading, before that distinction exists, would quietly
reinstate the association the later chapters worked hard to separate.
If the manuscript's actual Part III title still reads ``\ldots and
Obstructions,'' it should be changed to match.)

\begin{longtable}[]{@{}
  >{\raggedright\arraybackslash}p{(\columnwidth - 6\tabcolsep) * \real{0.2500}}
  >{\raggedright\arraybackslash}p{(\columnwidth - 6\tabcolsep) * \real{0.2500}}
  >{\raggedright\arraybackslash}p{(\columnwidth - 6\tabcolsep) * \real{0.2500}}
  >{\raggedright\arraybackslash}p{(\columnwidth - 6\tabcolsep) * \real{0.2500}}@{}}
\toprule\noalign{}
\begin{minipage}[b]{\linewidth}\raggedright
Correspondence
\end{minipage} & \begin{minipage}[b]{\linewidth}\raggedright
Tested in
\end{minipage} & \begin{minipage}[b]{\linewidth}\raggedright
Status
\end{minipage} & \begin{minipage}[b]{\linewidth}\raggedright
Qualification / open obligation
\end{minipage} \\
\midrule\noalign{}
\endhead
\bottomrule\noalign{}
\endlastfoot
Truth-value gaps ↔ underdetermination (absence of both polarities) & Ch.9 & Direct --- established & Gap is not a third truth value bolted on; it's what ``no answer yet'' already means once the two support questions are separated. \\
Truth-value gluts ↔ overdetermination (retained conflicting evidence) & Ch.10 & Direct --- established & Sharper than ``some propositions are both true and false'' --- contradictory evidence need not be deleted merely for conflicting. \\
FDE's four values ↔ independent \(\sigma^+,\sigma^-\) & Ch.9--10, sharpened again in Ch.20 & Direct --- established & Became the entry point into Part VI's sheaf machinery, not merely a tableaux gloss. \\
Disjunctive syllogism failure ↔ restricted propagation & Ch.11 & Direct --- established, technical centerpiece & DS's second step is shown to silently assume positive/negative support are mutually exclusive --- false exactly at overdetermined points. Not an ad hoc patch. \\
LPm / classical recapture ↔ repair biased toward consistency & Ch.12 & Direct --- established & Pays the cost of a glut only where evidence leaves no alternative; classical elsewhere. \\
FDE merge ↔ glut-free reconstruction obstruction & Ch.12, restated at cover level Ch.23 & Direct --- established, order-theoretic & Deliberately not cohomological in Part III; the cover-level generalization was checked for compatibility during Part VI's final pass and confirmed to be the same operation, not merely analogous (see Ch.23 row below). \\
Relevant logics (ternary relation) ↔ propagation licensing & Ch.11 gestures at it via the DS mechanism & Open obligation & Named missing result: expressive equivalence to the ternary relation, neither shown nor disproven. \\
\end{longtable}

\hypertarget{part-iv-construction-vagueness-and-degree}{%
\section{Part IV --- Construction, Vagueness, and Degree}\label{part-iv-construction-vagueness-and-degree}}

\begin{longtable}[]{@{}
  >{\raggedright\arraybackslash}p{(\columnwidth - 6\tabcolsep) * \real{0.2500}}
  >{\raggedright\arraybackslash}p{(\columnwidth - 6\tabcolsep) * \real{0.2500}}
  >{\raggedright\arraybackslash}p{(\columnwidth - 6\tabcolsep) * \real{0.2500}}
  >{\raggedright\arraybackslash}p{(\columnwidth - 6\tabcolsep) * \real{0.2500}}@{}}
\toprule\noalign{}
\begin{minipage}[b]{\linewidth}\raggedright
Correspondence
\end{minipage} & \begin{minipage}[b]{\linewidth}\raggedright
Tested in
\end{minipage} & \begin{minipage}[b]{\linewidth}\raggedright
Status
\end{minipage} & \begin{minipage}[b]{\linewidth}\raggedright
Qualification / open obligation
\end{minipage} \\
\midrule\noalign{}
\endhead
\bottomrule\noalign{}
\endlastfoot
Intuitionistic construction ↔ reconstruction path & Ch.13 & Direct --- established, sharpened & Pinned specifically to \(\neg\neg p \to p\) failing, not a looser ``existence from absurdity'' slogan --- the earlier looser version would have overstated the correspondence. \\
Support vs.~constructibility as one axis vs.~two & Ch.13 & Rejected as one axis; established as two orthogonal coordinates & A non-constructive derivation and a genuine construction can both count as positive support; Part III's apparatus doesn't (and didn't need to) distinguish them. \\
Many-valued/fuzzy degree ↔ resolution-stability & Ch.14 & Bridge, explicitly labeled as this book's added interpretation & Corrected from an earlier draft that stated this as though many-valued semantics already meant it. \(v_\lambda(p)\) vs.~the resolution profile \(\lambda \mapsto v_\lambda(p)\) now kept distinct. \\
Sorites / vagueness ↔ dispersion obstruction & Ch.15 & Bridge, developed & Corrected definition: accumulation of small variation, not exact-agreement-everywhere (\(\sum_i \epsilon_i\) need not be small even when each \(\epsilon_i\) is negligible). The earlier ``N and B both fail to fit \{T,F\}'' framing this traces back to is now understood to have been too coarse; superseded, not merely refined. \\
Dispersion obstruction ↔ Ch.24's proposed topology & Ch.15, tested directly against Ch.24 §3 during Part VI's final pass & \textbf{Distinct but related} & Ch.24's topology is qualitative-only (neighborhoods, no metric); dispersion needs the accumulation structure Ch.15 introduced. Confirmed by direct test, not assumed either way. \\
\end{longtable}

\hypertarget{part-v-identity-change-and-worldhood}{%
\section{Part V --- Identity, Change, and Worldhood}\label{part-v-identity-change-and-worldhood}}

\begin{longtable}[]{@{}
  >{\raggedright\arraybackslash}p{(\columnwidth - 6\tabcolsep) * \real{0.2500}}
  >{\raggedright\arraybackslash}p{(\columnwidth - 6\tabcolsep) * \real{0.2500}}
  >{\raggedright\arraybackslash}p{(\columnwidth - 6\tabcolsep) * \real{0.2500}}
  >{\raggedright\arraybackslash}p{(\columnwidth - 6\tabcolsep) * \real{0.2500}}@{}}
\toprule\noalign{}
\begin{minipage}[b]{\linewidth}\raggedright
Correspondence
\end{minipage} & \begin{minipage}[b]{\linewidth}\raggedright
Tested in
\end{minipage} & \begin{minipage}[b]{\linewidth}\raggedright
Status
\end{minipage} & \begin{minipage}[b]{\linewidth}\raggedright
Qualification / open obligation
\end{minipage} \\
\midrule\noalign{}
\endhead
\bottomrule\noalign{}
\endlastfoot
Necessary/contingent identity, rigid designation ↔ trajectory-first ontology & Ch.16 & Bridge & Explicitly not claimed to derive rigid designation --- rigid designation is the comparison case, not a target reduced to this book's account. \\
Identity ↔ invariant of admissible continuation & Ch.16 & Direct --- established, with one open dependency & The thesis itself is stated cleanly; but see the \(R_w\)/continuation row below --- this chapter's own notion of ``admissible'' is underspecified once regimes vary by state. \\
Priest's contradictory transition instant ↔ continuation-interface object & Ch.17 & Tension --- retained, as a productive alternative & Explicitly qualified ``sometimes, not always'' --- must not become a general argument against dialetheism, given Part III's four chapters on genuine gluts. Not a derivation of Priest's account; a rival with a named scope limit. \\
Hegelean motion ↔ RSVP field dynamics & Ch.18 & Tension --- retained & Not upgraded despite RSVP doing real argumentative work; the ontological-fact-vs-inadequate-vocabulary question is left open on purpose. \\
Leibniz Continuity Condition ↔ discrete repair events & Ch.19 & Tension --- retained, given a specific test & ``Does a continuous trajectory underlie every discrete repair at some finer description'' is posed as a real either/or, not assumed to reconcile. \\
Spread Hypothesis / temporal flow ↔ Chain of Memory & Ch.19 & Bridge, deliberately narrowed & Explicitly not a theory of physical time --- narrowed to instantaneous-state ≠ history-indexed-state. Overselling this was an identified risk and avoided. \\
\textbf{\(R_w\) (Ch.7) ↔ admissible continuation (Ch.16)} & Ch.16, surfaced during Part VI's final pass & \textbf{New open research problem --- not resolved during reconciliation, on purpose} & \(x_{t} \xrightarrow{?} x_{t+1}\) with \(R_{x_t} \neq R_{x_{t+1}}\): ``admissible continuation'' doesn't yet say relative to which regime a transition step is judged --- \(R_{x_t}\), \(R_{x_{t+1}}\), an intersection/union, a transport of one into the other, or a distinct transition regime \(R_\gamma\). Ch.17's transition-as-its-own-type move suggests the last option but was correctly not promoted to a definition. This is a genuine result of completing all six parts, not a loose end to patch quietly. \\
\end{longtable}

\hypertarget{part-vi-local-logic-and-global-obstruction}{%
\section{Part VI --- Local Logic and Global Obstruction}\label{part-vi-local-logic-and-global-obstruction}}

\begin{longtable}[]{@{}
  >{\raggedright\arraybackslash}p{(\columnwidth - 6\tabcolsep) * \real{0.2500}}
  >{\raggedright\arraybackslash}p{(\columnwidth - 6\tabcolsep) * \real{0.2500}}
  >{\raggedright\arraybackslash}p{(\columnwidth - 6\tabcolsep) * \real{0.2500}}
  >{\raggedright\arraybackslash}p{(\columnwidth - 6\tabcolsep) * \real{0.2500}}@{}}
\toprule\noalign{}
\begin{minipage}[b]{\linewidth}\raggedright
Correspondence
\end{minipage} & \begin{minipage}[b]{\linewidth}\raggedright
Tested in
\end{minipage} & \begin{minipage}[b]{\linewidth}\raggedright
Status
\end{minipage} & \begin{minipage}[b]{\linewidth}\raggedright
Qualification / open obligation
\end{minipage} \\
\midrule\noalign{}
\endhead
\bottomrule\noalign{}
\endlastfoot
Local valuations ↔ local admissible sections \(s_i \in \mathcal{A}(U_i)\) & Ch.20 & Direct --- established & \\
Compatibility of local descriptions ↔ Čech coboundary vanishing & Ch.21 & Direct --- established, corrected apparatus & Not literal \(H^1\): for any chosen 0-cochain, \(\delta s\) is a coboundary by construction, so its class in \(H^1\) is trivially zero regardless of agreement. Renamed \textbf{coherence defect} \(\delta^\pm(p) \in \check{C}^1\), a cochain-level check. \\
Failure of global coherent determination ↔ nonzero \(H^1\) / glut & Ch.20--22 & \textbf{Rejected as originally stated}, replaced by the Mutual Non-Determination Proposition & ``Glut \(= [g] \neq 0\)'' is not a valid identification --- glut is semantic (dual support), \(H^1\)-type obstruction is geometric. The corrected result: coherence status and the FDE merge \(\Sigma(p)\) determine neither each other, proven by an explicit six-cell realizability table (Ch.22), not asserted by analogy. \\
Bilattice join ↔ ordinary sheaf gluing & Ch.20--21 & \textbf{Rejected} & Join is an information-preserving merge/repair of a non-matching family; gluing requires a family that already agrees on overlaps. Different operations under the same informal word ``combine.'' Terminology fixed going forward: restriction ≠ gluing ≠ repair. \\
TARTAN ↔ local-to-global logical regime assembly & Ch.24 §3 (proposed) & Open obligation & Two named missing results: does refinement only reveal genuine coherence defects or manufacture spurious ones; is the resulting limit/colimit well-behaved. \\
Ch.12's merge \(M\) ↔ Ch.23's merge \(\Sigma\) & Ch.23, checked directly during Part VI's final pass & Direct --- identified & \(\Sigma = M\) under cover indexing; the name change marks the structural generalization (bare sources → cover-indexed patches), not a different operation. \\
\(T,F,B,N\) ↔ fibers of a projection, not points of a richer space & Ch.24 §1 & Direct --- established & A restatement of the Mutual Non-Determination Proposition in geometric language, not new content --- the only fully rigorous claim in Ch.24. \\
Repair cost as a metric on \(\mathcal{S}(p)\) & Ch.24 §4 (proposed) & Open obligation & Two named missing results: whether an action-independent repair distance exists, and whether it satisfies a triangle inequality. \\
\end{longtable}

\begin{center}\rule{0.5\linewidth}{0.5pt}\end{center}

\begin{center}\rule{0.5\linewidth}{0.5pt}\end{center}

\hypertarget{unpursued-correspondences}{%
\section{Unpursued correspondences}\label{unpursued-correspondences}}

Two entries inherited from the original planning-stage matrix never had
a chapter written for them. They are not weak Bridges or exposed
obligations --- the manuscript contains no attempt to test them at all,
and listing them in the main tables above would misstate what this
ledger claims to record (status \emph{after drafting}, for material the
finished six parts actually contain).

\begin{longtable}[]{@{}
  >{\raggedright\arraybackslash}p{(\columnwidth - 4\tabcolsep) * \real{0.3333}}
  >{\raggedright\arraybackslash}p{(\columnwidth - 4\tabcolsep) * \real{0.3333}}
  >{\raggedright\arraybackslash}p{(\columnwidth - 4\tabcolsep) * \real{0.3333}}@{}}
\toprule\noalign{}
\begin{minipage}[b]{\linewidth}\raggedright
Correspondence
\end{minipage} & \begin{minipage}[b]{\linewidth}\raggedright
Original proposal
\end{minipage} & \begin{minipage}[b]{\linewidth}\raggedright
Why it stayed unpursued
\end{minipage} \\
\midrule\noalign{}
\endhead
\bottomrule\noalign{}
\endlastfoot
Priest's semantic closure (self-referential truth predicate) ↔ Spherepop's history-based operational semantics & A self-referential object language and a computation model that can refer to its own prior states looked structurally similar & No chapter in Parts I--VI was assigned this topic; the six-part outline never routed through it \\
Priest's naive/paraconsistent set theory (material, relevant, and model-theoretic strategies) ↔ obstruction as candidate propagation rules & Priest's three strategies for handling set-theoretic paradoxes looked like candidate instances of this book's propagation-restriction machinery & Same reason --- no chapter's scope included it \\
\end{longtable}

This isn't a defect in the finished book. A twenty-four-chapter argument
built around one central hierarchy doesn't need to absorb every topic
Priest covers across two books, and both of these are now more
approachable than they would have been before this project's framework
existed --- they're candidates for separate future work precisely because
the machinery to attempt them properly (the support/coherence/repair
apparatus, the admissibility regime vocabulary) now exists in a form
mature enough not to require inventing it from scratch alongside the
comparison itself.

\begin{center}\rule{0.5\linewidth}{0.5pt}\end{center}

\hypertarget{what-changed-as-a-record-rather-than-a-silent-overwrite}{%
\section{What changed, as a record rather than a silent overwrite}\label{what-changed-as-a-record-rather-than-a-silent-overwrite}}

Three correspondences were investigated and found to conflate distinct
structures, and are marked Rejected rather than quietly dropped: glut as
a nonzero cohomology class, bilattice join as ordinary sheaf gluing, and
non-normal worlds as refusal. All three were plausible enough to have
been rated Direct or treated as settled before the relevant chapter was
actually drafted. That they didn't survive contact with the chapters is
evidence the drafting process constrained its own conclusions rather than
rationalizing whatever correspondence was proposed first.

One new open research problem --- the \(R_w\)/continuation transition rule ---
was generated by completing Parts II and V, not resolved by their
completion. It's recorded here deliberately unresolved, per instruction,
since fixing it during this reconciliation pass would erase evidence that
finishing the book found a real gap rather than merely filling in a
pre-drawn outline.

This ledger is now the right basis for the book's introduction and
conclusion, and for identifying which individual results (the Mutual
Non-Determination Proposition and its realizability proof, most likely,
along with the \(R_w\)/continuation problem) might warrant a standalone
paper ahead of the full manuscript.

\textbf{This revision's specific change}: split the undifferentiated Bridge
category into Bridge (diffuse gap) and Open obligation (one nameable
missing result), reclassifying five rows accordingly, and moved the two
correspondences the manuscript never actually tested (semantic
closure/Spherepop, naive set theory) out of the main tables into their
own section. The ledger's claim to record what the finished manuscript
established was slightly inconsistent with carrying two rows the
manuscript never attempted; that inconsistency is now fixed rather than
carried forward again.
